\chapter{Validation}\label{ch:validation}

This chapter checks the implementation against the requirements of
Chapter~\ref{ch:analysis} and the protocol described in Chapter~\ref{ch:project}. Because
the protocol's own correctness criterion is exact, a run either converges to \emph{the}
unique minimum spanning tree of the current graph or it does not, validation is built
around a single reference check applied uniformly to every scenario, rather than a
collection of hand-picked expected outputs. What differs from test to test is only how
that scenario is produced: which topology, which sequence of events, and which of the two
runtimes executes it.

\section{The correctness oracle}\label{sec:oracle}

Every test in the suite ultimately delegates its verification to a single call:
\texttt{sim/oracle.check(graph, summaries)}. This function validates a converged execution
against \texttt{graph.kruskal} and accepts the run if and only if all four of the
following hold:

\begin{itemize}
  \item \textbf{Quiescence:} Every node reports \texttt{halted: True}, guaranteeing that
  the protocol reached stability;
  \item \textbf{Edge Set:} The union of all locally reported tree edges is identical to
  the exact edge set produced by \texttt{kruskal};
  \item \textbf{Root Count:} The number of nodes reporting \texttt{parent: None} equals
  the exact number of connected components in the current graph, ensuring precisely one
  root per fragment;
  \item \textbf{Structural Integrity:} Every non-root node's parent pointer connects to an
  active tree edge, confirming that directed parent references form a valid rooted forest.
\end{itemize}

Because the total edge ordering resolves all weight ties, the reference MST is strictly
unique: there is exactly one valid topology and edge set that the oracle will accept.

\section{Two execution paths, one protocol}

Every test drives one of two runtimes: the real OTP
actor system, or \texttt{test/sim/runner.gleam}'s process-free simulator. Both call exactly
the same \texttt{d2mst/algorithm.handle} on exactly the same protocol modules; the only
difference is what delivers events to it. The simulator interprets a node's \texttt{Send}
effects directly into its own FIFO queue instead of onto a real link process, and lets a
scenario be driven completely deterministically: the same sequence of \texttt{wake},
\texttt{fail\_link}, \texttt{add\_link}, and \texttt{crash\_node} calls produces the exact
same interleaving of events every time, and can be stepped one message at a time with
\texttt{step\_one}.

The actor runtime is not redundant with this. It is the only place where a
\texttt{LinkDown} is an actual monitored BEAM process dying, rather than an event handed
directly to the queue, and the only place where two nodes' handlers genuinely run
concurrently instead of being serialized by a single driving loop. Section~\ref{sec:actor}
describes what it is used to check that the simulator structurally cannot.

\section{Randomized and regression testing}\label{sec:fuzz}

Fixed graphs catch bugs that are wrong for every input; they are not good for bugs that
depend on a particular tree shape or a particular pair of overlapping events. To avoid this 
\texttt{engine/generator.gleam} implements a seeded connected-graph builder: the same seed 
always produces the same graph, so a fuzz failure is reproducible from its integer seed with 
no additional state to capture. Most fuzz tests follow the same shape: converge on a batch 
of seeded topologies, apply a scenario, and check the oracle after every one.

A second group of tests fires several topology events \emph{without settling in between},
so that their protocol rounds are genuinely concurrent and not just sequential with no gap.

The most targeted tests are fixed, hand-built
graphs written to reconstruct a specific race the general-purpose fuzzing above did not
happen to find, and are kept as permanent regression tests once the underlying bug is
fixed. For example \texttt{overlapping\_cycles\_test} builds two link additions whose cycles genuinely
share tree edges, exactly the race Chapter~\ref{ch:project}'s "overlapping cycles
serialization" paragraph describes.

\section{Long-running randomized simulation}\label{sec:long-running}

Chapter~\ref{ch:intro}'s system overview commits the evaluation to, among other scenarios,
"concurrent random additions and deletions" and "long-running simulations with node and
edge events occurring at different rates". \texttt{long\_running\_random\_topology\_test}
is built directly for this: starting from a five-hundred-node random graph, it fires sixty
topology events, drawn uniformly across failing a random edge, crashing a random node,
adding a random edge, and joining a new isolated node, biased against crashing below five
surviving nodes so the run always has something left to mutate, in bursts of five with
\emph{nothing allowed to settle in between}. The oracle of Section~\ref{sec:oracle} is
checked after every burst, so a burst of five overlapping events
must both leave the system in a checkable state and eventually converge correctly, not just
the final one. \texttt{fuzz\_long\_running\_random\_topology\_test} repeats the same
procedure from five different starting topologies. 

\section{Interactive demonstration}

\texttt{gleam run} (Appendix~\ref{app:devenv}) is not part of the automated suite, but
complements it: it builds a random connected graph, brings it to a converged MST, and then
repeatedly applies
a whole batch of node and edge failures and additions as one overlapping burst, waits for
convergence, checks the result against \texttt{graph.kruskal}, and reports the message
count and wall-clock time for the round.

\section{Message-complexity benchmark}\label{sec:bench}

Chapter~\ref{ch:project} derives $O(|E|)$ message complexity for the naive Phase 3 search
and $O(N \log N)$ for the binary-search alternative. \texttt{test/bench\_report.gleam} is
not a \texttt{gleeunit} test, it can be run with:

{\footnotesize
\begin{verbatim}
gleam build --target erlang
erl -pa build/dev/erlang/*/ebin -noshell \
  -eval "bench_report:run(), halt()."
\end{verbatim}
}

For each graph size $N \in \{10, 20, 40, 80, 160, 320\}$ and each burst size $k \in \{1, 2,
4, 8, 16\}$ (how many tree edges, or non-adjacent pairs, are perturbed at once before the
network is allowed to settle), the benchmark builds \texttt{seed\_count} $= 15$ random
graphs, runs the scenario, measures the messages sent between before and after the burst,
and reports the mean alongside its ratio to the complexity claim's asymptotic term: mean
$|E|$ for the naive search, $N \log_2 N$ for the binary search, and $N$ for the addition
protocol. Failure scenarios are measured at two densities, a sparse graph
($\mathrm{edge\_pct} = 15$) and a dense one ($\mathrm{edge\_pct} = 85$), since the naive
bound is stated in terms of $|E|$ and the two densities are where that quantity differs
most from $N$.
The full sweep, one table per scenario:

{\footnotesize
\begin{verbatim}
Failure - naive Phase 3, sparse (edge_pct=15)
N              k=1           k=2            k=4            k=8            k=16
--------------------------------------------------------------------------------
N=10     111.0 (7.9x)  132.8 (9.5x)  150.7 (10.8x)  106.4 (7.6x)              n/a
N=20     334.7 (7.3x) 477.1 (10.5x)  683.8 (15.0x)  863.0 (18.9x)  899.4 (19.7x)
N=40    1105.0 (7.4x) 1549.5 (10.3x) 2436.1 (16.3x) 3471.3 (23.2x) 4795.3 (32.0x)
N=80    3941.5 (7.3x) 5580.8 (10.3x) 8483.7 (15.7x) 12171.3 (22.5x) 17930.3 (33.1x)
N=160  13979.1 (6.8x) 20334.7 (9.9x) 29249.5 (14.3x) 46246.5 (22.6x) 68519.5 (33.4x)
N=320  53579.5 (6.7x) 75717.4 (9.5x) 112743.7 (14.1x) 172983.5 (21.7x) 258182.9 (32.4x)

Failure - naive Phase 3, dense (edge_pct=85)
N              k=1           k=2            k=4            k=8            k=16
--------------------------------------------------------------------------------
N=10     265.6 (6.7x)  333.8 (8.4x)  399.3 (10.0x)  427.1 (10.8x)             n/a
N=20    1011.6 (6.2x) 1229.9 (7.5x)  1565.5 (9.6x) 1926.9 (11.8x) 2140.8 (13.1x)
N=40    3980.8 (6.0x) 4922.5 (7.4x)  6325.5 (9.5x) 8036.7 (12.1x) 9816.5 (14.7x)
N=80   14937.9 (5.5x) 18422.4 (6.8x) 22465.1 (8.3x) 28189.5 (10.4x) 36704.1 (13.6x)
N=160  64180.2 (5.9x) 73621.7 (6.8x) 93659.2 (8.6x) 110659.1 (10.2x) 143435.5 (13.2x)
N=320 241979.6 (5.6x) 311050.5 (7.2x) 368015.0 (8.5x) 456982.8 (10.5x) 577752.4 (13.3x)

Failure - binary search Phase 3, sparse (edge_pct=15)
N              k=1           k=2            k=4            k=8            k=16
--------------------------------------------------------------------------------
N=10     119.6 (3.6x)  151.5 (4.6x)  173.3 (5.2x)  117.6 (3.5x)              n/a
N=20     345.9 (4.0x)  571.8 (6.6x)  979.9 (11.3x) 1428.3 (16.5x) 1635.6 (18.9x)
N=40    1044.3 (4.9x) 1793.1 (8.4x) 3258.5 (15.3x) 5372.4 (25.2x) 8383.4 (39.4x)
N=80    2587.9 (5.1x) 4742.0 (9.4x) 8891.1 (17.6x) 15617.3 (30.9x) 26616.7 (52.6x)
N=160   5699.9 (4.9x) 10948.5 (9.3x) 21327.9 (18.2x) 40897.2 (34.9x) 73864.5 (63.1x)
N=320  13499.3 (5.1x) 26275.0 (9.9x) 50914.7 (19.1x) 97952.9 (36.8x) 186660.3 (70.1x)

Failure - binary search Phase 3, dense (edge_pct=85)
N              k=1           k=2            k=4            k=8            k=16
--------------------------------------------------------------------------------
N=10     187.8 (5.7x)  297.0 (8.9x)  424.6 (12.8x)  479.5 (14.4x)             n/a
N=20     518.7 (6.0x) 881.5 (10.2x) 1545.9 (17.9x) 2232.6 (25.8x) 2806.1 (32.5x)
N=40    1326.1 (6.2x) 2412.9 (11.3x) 4407.7 (20.7x) 7547.1 (35.5x) 11555.8 (54.3x)
N=80    3178.1 (6.3x) 5700.4 (11.3x) 10558.1 (20.9x) 18427.5 (36.4x) 31333.2 (62.0x)
N=160   6882.9 (5.9x) 13367.6 (11.4x) 26449.0 (22.6x) 49068.5 (41.9x) 89685.5 (76.6x)
N=320  15734.3 (5.9x) 30885.2 (11.6x) 59553.6 (22.4x) 114585.3 (43.0x) 220605.3 (82.8x)

Addition, edge_pct=40
N              k=1           k=2            k=4            k=8            k=16
--------------------------------------------------------------------------------
N=10      26.7 (2.7x)  55.2 (5.5x)  111.1 (11.1x)  217.7 (21.8x)   438.0 (43.8x)
N=20      40.6 (2.0x)  84.4 (4.2x)   165.1 (8.3x)  333.6 (16.7x)   678.1 (33.9x)
N=40      76.8 (1.9x) 147.5 (3.7x)   297.5 (7.4x)  595.9 (14.9x)  1194.7 (29.9x)
N=80     129.3 (1.6x) 251.7 (3.1x)   500.5 (6.3x)  994.1 (12.4x)  2041.7 (25.5x)
N=160    227.7 (1.4x) 454.0 (2.8x)   894.3 (5.6x) 1774.6 (11.1x)  3534.5 (22.1x)
N=320    410.1 (1.3x) 812.1 (2.5x)  1613.9 (5.0x) 3254.2 (10.2x)  6452.5 (20.2x)
\end{verbatim}
}

At the lowest concurrency, $k=1$, a single isolated event, which is what
Chapter~\ref{ch:project}'s complexity derivations actually analyze, both failure tables'
ratio columns stay close to a constant as $N$ grows by a factor of thirty-two, from $N=10$
to $N=320$ (naive: roughly $5.5$-$8\times$ in both densities; binary search: roughly
$4$–$6.3\times$, excluding the $N=10$ sparse outlier discussed below), which is exactly what
an $O(|E|)$ and an $O(N \log N)$ bound respectively predict. The naive/binary-search
comparison also confirms the motivation Chapter~\ref{ch:project} gives for the binary-search
procedure in the first place: at $N=160$, binary search sends $5700$ messages against the
naive search's $13979$ in the sparse graph ($2.5\times$ fewer), and $6883$ against
$64180$ in the dense one ($9.3\times$ fewer). The gap widens
with $N$ too: at $N=320$ binary search is $4.0\times$ fewer messages in the sparse graph and
$15.4\times$ fewer in the dense one, both wider than at $N=160$. At the
smallest size, $N=10$, the two are close and binary search is even marginally more expensive
in the sparse case ($119.6$ against $111.0$): the extra broadcast/convergecast rounds that
selecting a median a few times costs are not yet paid back by a saving on $|E|$ that is
still tiny at that size, so the asymptotically better procedure only earns its keep once the
graph is large enough.

The ratio columns tell us something different. For naive, the $k=8$ ratio climbs with $N$ only up to a point and then flattens:
$7.6\times$, $18.9\times$, $23.2\times$, $22.5\times$, $22.6\times$, $21.7\times$ across
$N=10,\dots,320$ in the sparse graph, and just $10.8\times$-$13.6\times$ across the same
range in the dense one; the $k=16$ column repeats the same shape one notch higher, settling
around $32$-$33\times$ sparse and $13$-$15\times$ dense from $N=80$ onward. Binary search
never settles: its $k=8$ ratio keeps climbing across the whole range, from $3.5\times$ to
$36.8\times$ sparse and $14.4\times$ to $43.0\times$ dense, and $k=16$ climbs even more
steeply, reaching $70.1\times$ sparse and $82.8\times$ dense at $N=320$. Both bounds are derived for a single, isolated repair event, so this divergence is itself
evidence that concurrency affects the two strategies differently. Firing $k$ failures without letting
the system settle in between means a fragment identity change from one repair can hit a
Phase 3 search already in flight for another, and the implementation restarts such a search
from scratch.

Naive's Phase 3 search completes in one broadcast-down/convergecast-up pass, so it is
exposed to this risk only briefly, and a restart repeats work of the same order as the
search itself, $O(|E|)$. Binary search instead narrows its candidate interval across
several sequential broadcast/convergecast rounds, and a restart
discards that narrowed $[a,b]$ interval along with the sample built up over it, forcing the
search to begin again from an unfiltered scan. A search that stays in flight for more
rounds is exposed to invalidation for longer, and larger fragments need more such round,
which would explain why the binary-search ratio keeps growing with $N$ where naive's does
not.

The addition table shows the opposite trend under concurrency: at $k=8$ the ratio to $N$
\emph{falls} as $N$ grows, from $21.8\times$ at $N=10$ to $10.2\times$ at $N=320$, at
$k=16$ it falls from $43.8\times$ to $20.2\times$, and at $k=1$ it falls from $2.7\times$ to
$1.3\times$, comfortably inside the $O(N)$ bound at every cell. This matches the addition
protocol's different answer to concurrency: rather than letting an in-progress search be
invalidated and restarted, same-fragment additions are serialized one at a time behind the
root's \msg{Privilege} token (Chapter~\ref{ch:project}'s "overlapping cycles
serialization"), so $k$ concurrent additions cost closer to $k$ times one addition's cost
plus a fixed coordination overhead, rather than $k$ searches each separately paying a
restart penalty that grows with how large a fraction of the fragment's search state a
restart invalidates.

\section{Coverage against the requirements of Chapter~\ref{ch:analysis}}

Fault tolerance and eventual convergence are what the oracle of Section~\ref{sec:oracle}
checks directly, on every scenario in this chapter; partition tolerance is exercised by
every test in which a crash or failure actually disconnects the graph and the oracle still
requires exactly one root per surviving component. Failure transparency follows from
Section~\ref{sec:actor}'s process-death tests: a node never has to distinguish a crashed
neighbor from a plain link failure, because \texttt{link.gleam} already reduced one to the
other before the protocol saw an event. Concurrency transparency is the specific target of
Section~\ref{sec:fuzz}'s overlapping-event tests and Section~\ref{sec:long-running}.
